\documentclass[11pt,a4paper]{article}

\usepackage[utf8]{inputenc}
\usepackage[T1]{fontenc}
\usepackage[turkish]{babel}
\usepackage[margin=2.4cm]{geometry}
\setlength{\headheight}{14pt}
\usepackage{graphicx}
\usepackage{float}
\usepackage{booktabs}
\usepackage{tabularx}
\usepackage{amsmath,amssymb}
\usepackage{gensymb}
\usepackage[hidelinks]{hyperref}
\usepackage{fancyhdr}
\usepackage{enumitem}
\usepackage{xcolor}

\definecolor{demoblue}{rgb}{0.10,0.24,0.62}
\definecolor{armgreen}{rgb}{0.10,0.44,0.20}
\definecolor{thesisred}{rgb}{0.62,0.15,0.15}

\pagestyle{fancy}
\fancyhf{}
\fancyhead[L]{\small Hafta 4 -- Demo Sistemi, Ust Ekstremite ve Tez Taslagi}
\fancyhead[R]{\small \thepage}
\renewcommand{\headrulewidth}{0.4pt}

\begin{document}

\begin{center}
    {\LARGE\bfseries Hafta 4 Gelistirme Raporu:\\[4pt]
    \textcolor{demoblue}{Ornek Hareketleri Gosterme Sistemi},\\[4pt]
    \textcolor{armgreen}{Daha Kararli Kol Takibi} ve\\[4pt]
    \textcolor{thesisred}{Bitirme Tezi Taslagi}}\\[12pt]
    {\normalsize Nisan 2026}
\end{center}

\vspace{6pt}

\begin{abstract}
Bu rapor, proje gelisim surecinin dorduncu haftasinda tamamlanan uc ana
calismayi belgelemektedir. Ilk olarak, oyunlastirilmis tarama gorevleri icin
kullaniciya ornek hareketleri gosteren \textbf{ghost avatar tabanli demo sistemi}
tasarlanmis ve gorev tanimlariyla butunlestirilmistir. Ikinci olarak,
ust ekstremite izleme hattinda \textbf{omuz referans cercevesi},
\textbf{kalibrasyon tabanli goreli rotasyonlar} ve \textbf{avatar yansitma}
katmanlari tanimlanarak kol takibinin daha stabil ve tekrar edilebilir hale
gelmesi saglanmistir. Ucuncu olarak, projenin akademik teslimi icin
\textbf{bitirme tezi taslagi} universite formatina uygun bicimde yazilmistir.
Bu uc ciktinin ortak amaci; sistemin teknik olgunlugunu, kullanici
yonlendirme kapasitesini ve akademik raporlanabilirligini ayni anda ileri
tasimaktir.
\end{abstract}

\tableofcontents
\vspace{8pt}

% ================================================================
\section{Rapor Kapsami ve Hafta 3 Uzerine Eklenenler}
\label{sec:kapsam}

Hafta~3 raporu agirlikli olarak full-body IK altyapisinin gercek zamanli
uygulamasina, shin-tracker ve 5-tracker modlarina, omurga cozumune,
editor araclarina ve sahne ortam tasarimina odaklanmisti. Hafta~4'te ise
odak, ayni altyapiyi \textbf{kullaniciya gosteren}, \textbf{daha kararli}
hale getiren ve \textbf{akademik olarak raporlayan} ust katmanlara kaymistir.

\begin{table}[H]
\centering
\caption{Hafta 3 ve Hafta 4 odak farki}
\label{tab:hafta34}
\renewcommand{\arraystretch}{1.25}
\begin{tabularx}{\textwidth}{lXX}
\toprule
\textbf{Baslik} & \textbf{Hafta 3} & \textbf{Hafta 4} \\
\midrule
Ana hedef & Kinematik altyapiyi uygulamak & Uygulanan sistemi gostermek, stabil hale getirmek ve tezlestirmek \\
Kullanici deneyimi & IK ve cevre kurulumu & Ghost avatar ile ornek hareket gosterimi \\
Ust ekstremite & Temel IK/FK genislemesi & Referans cercevesi, goreli rotasyon, pose follower ve retargeting \\
Dokumantasyon & Teknik implementasyon raporu & Haftalik rapor + bitirme tezi taslagi \\
\bottomrule
\end{tabularx}
\end{table}

Bu rapor, hafta~4'te tamamlanan uc somut teslim etrafinda yapilandirilmistir:
\begin{enumerate}[nosep]
    \item Gorev bazli \textbf{ornek hareketleri gosteren demo sistemi},
    \item Kol track islemini daha stabil hale getiren \textbf{ust ekstremite izleme/avatar} katmani,
    \item Projenin akademik teslimi icin yazilan \textbf{bitirme tezi taslagi}.
\end{enumerate}

% ================================================================
\section{Ornek Hareketleri Gosteren Sistem}
\label{sec:demo}

Oyunlastirilmis tarama uygulamalarinda yazili yonlendirme tek basina yeterli
degildir. Ozellikle mini squat, kontrollu inis, tek ayak squat ve reach
gorevleri gibi hareketlerde kullanicinin neyi ne kadar yapmasi gerektigi,
bir \textit{gorsel referans} olmadan farkli sekillerde yorumlanabilir.
Bu nedenle hafta~4'te, gorev aciklamalarina paralel calisan bir
\textbf{ghost avatar demo sistemi} hazirlanmistir.

\subsection{ScriptableObject Tabanli Poz Veri Modeli}

Sistem, her gorev icin poz bilgisini sahne icine hard-code etmek yerine
asset tabanli bir veri modeli uzerine kurulmustur. Bu sayede yeni gorev
eklemek veya mevcut gosteri hareketlerini guncellemek icin kod degistirmek
gerekmez; ilgili pose asset'i guncellemek yeterlidir.

\begin{table}[H]
\centering
\caption{Demo sisteminin temel bilesenleri}
\label{tab:demo_bilesen}
\renewcommand{\arraystretch}{1.22}
\begin{tabularx}{\textwidth}{lX}
\toprule
\textbf{Bilesen} & \textbf{Sorumluluk} \\
\midrule
\texttt{PoseSnapshotSO} & Tek bir keyframe'in kalca pozisyonu ve tum eklem rotasyonlarini saklar \\
\texttt{PoseSequenceSO} & Bir goreve ait sirali keyframe listesini ve gecis hizini tutar \\
\texttt{TaskDefinition.demoSequence} & Her gorevin hangi demo sekansi ile eslesecegini belirler \\
\texttt{PoseDemoController} & Ghost avatar uzerinde bu keyframe'leri oynatir \\
\texttt{GamificationSetup} & Demo asset'lerini ve sahne baglantilarini editor ortaminda otomatik kurar \\
\texttt{SessionConfigSO} & Gorevlerin klinik oncelige gore sirasini ve oturum akisini tanimlar \\
\bottomrule
\end{tabularx}
\end{table}

\texttt{PoseSnapshotSO}, hips local position verisini de tuttugu icin sadece
eklem rotasyonlari degil; comelme, alcak inis ve reach gibi durumlarda
gorulen \textbf{gobde-agirlik merkezi dususu} da kaydedilebilmektedir.
Bu, onceki sabit Euler yaklasimina gore daha gercekci bir gosterim sunar.

\subsection{Pozlar Arasi Gecis Mantigi}

Ghost avatar tek bir sabit poza atlamaz; keyframe'ler arasi yumusak gecis
uygular. Rotasyonlar kuaternion uzayinda \textit{Slerp}, kalca pozisyonu ise
dogrusal enterpolasyon ile guncellenir:

\begin{align}
    Q(t) &= \mathrm{Slerp}(Q_i, Q_{i+1}, \alpha)
    \label{eq:slerp_demo} \\
    \mathbf{p}_{\text{hip}}(t) &= (1-\alpha)\,\mathbf{p}_i + \alpha\,\mathbf{p}_{i+1}
    \label{eq:lerp_demo}
\end{align}

Burada $\alpha \in [0,1]$ gecis ilerlemesini, $Q_i$ ve $Q_{i+1}$ ise ardisik
keyframe rotasyonlarini temsil eder. Boylece gosterim sistemi ``fotograf
degitirir'' gibi degil, \textbf{hareket eder} gibi algilanir.

\subsection{Varsayilan Gorev Kutuphanesi}

\texttt{PoseDemoController} icinde, gorev tipine gore cagrilabilen bir
varsayilan demo kutuphanesi de tanimlanmistir. Bu kutuphane, ozellikle
asset atamasi yapilmamis gorevlerde sistemin bos kalmasini onler.
Varsayilan kutuphane asagidaki gorevleri kapsar:

\begin{itemize}[nosep]
    \item Dik durus,
    \item Saga egilme, sola egilme ve one egilme,
    \item Mini squat,
    \item Sag/sol tek ayak denge,
    \item Yerinde yurume simulasyonu,
    \item Kontrollu inis taramasi,
    \item Modifiye anterior Y-balance sag/sol,
    \item Tek ayak squat sag/sol.
\end{itemize}

Boylece yeni bir sahnede demo asset'leri henuz elle uretilmemis olsa bile,
sistem gorevleri gorsel olarak gosterebilecek minimum bir davranisa sahip
olur.

\subsection{Editor Otomasyonu ve Ghost Avatar Uretimi}

Hafta~4'te sadece runtime oynatici degil, ayni zamanda onu uretmeyi ve
baglamayi hizlandiran bir editor kurulumu da eklenmistir.
\texttt{GamificationSetup} araci tek tikla:

\begin{enumerate}[nosep]
    \item \texttt{TaskDefinition} asset'lerini olusturur veya yukler,
    \item Ghost avatar'i ana avatarin kopyasi olarak sahneye ekler,
    \item Demo keyframe ve sequence asset'lerini otomatik uretir,
    \item Klinik onerili gorev sirasi icin bir \texttt{SessionConfigSO} hazirlar,
    \item UI, skor ve gorev yonetimi bilesenlerini birbirine baglar.
\end{enumerate}

Ozellikle \texttt{CreateOrUpdateAutoDemoSequences()} akisinda,
varsayilan demo kutuphanesindeki pozlar ghost avatar'a uygulanmakta,
ardindan her biri birer \texttt{PoseSnapshotSO} asset'i olarak yakalanmakta
ve ilgili \texttt{PoseSequenceSO} icine yazilmaktadir. Bu yontem, editorde
olusturulan gosterim verisinin tekrar kullanilabilir ve surdurulebilir
olmasini saglar.

\subsection{Klinik Oturum Sirasi ile Butunlesme}

Demo sistemi yalnizca gorsel bir yardimci degil, ayni zamanda gorev akisini
standardize eden oturum yapisinin parcasi haline getirilmistir.
\texttt{SessionConfigSO.SortTasksByClinicalRecommendation()} ile gorevler,
klinik oncelige gore asagidaki gibi dizilmektedir:

\begin{center}
\textit{Standing $\rightarrow$ LandingScreen $\rightarrow$ MiniSquat
$\rightarrow$ SingleLegSquat $\rightarrow$ ModifiedYBalanceAnterior
$\rightarrow$ Lean gorevleri $\rightarrow$ SingleLegBalance
$\rightarrow$ WalkSimulation}
\end{center}

Bu siralama, daha fazla diz fleksiyonu ve valgus gozlemi gerektiren
gorevleri oturumun erken asamasina alarak tarama kalitesini artirmayi
hedeflemektedir.

% ================================================================
\section{Kol Takibini Daha Stabil Hale Getiren Ust Ekstremite Sistemi}
\label{sec:arm}

Hafta~4'te ust ekstremite calismalari, ``kol takip edilsin'' seviyesinden
``kol takibi daha kararli, yorumlanabilir ve avatara aktarilabilir
olsun'' seviyesine tasinmistir. Bu kapsamdaki ana katmanlar:
\texttt{UpperLimbKinematicsUI}, \texttt{IUpperLimbKinematicsSource},
\texttt{UpperLimbAvatarPoseFollower}, \texttt{UpperLimbAvatarRetargeter}
ve donanimsiz test icin \texttt{UpperLimbKinematicsTestSource}'dur.

\subsection{Omuz Referans Cercevesinin Secimi}

Kol izleme sistemlerinde en onemli sorunlardan biri, omuz hareketlerinin hangi
referans cercevesine gore yorumlanacagidir. Sadece dunya eksenlerine gore
olculen ust kol rotasyonlari, kullanici govdesi dondugunde veya HMD yonu
degistiginde kararsiz gorunebilir. Bu nedenle sistemde omuz referansi uc kademeli
olarak secilir:

\begin{enumerate}[nosep]
    \item Torso tracker varsa ve takipteyse: \textbf{birincil referans},
    \item Torso tracker yoksa: HMD veya alternatif omuz referansi,
    \item Ikisi de yoksa: dunya eksenleri.
\end{enumerate}

Bu secim, \texttt{ResolveShoulderFrame()} mantigi ile yapilir. Sonuc olarak
omuz abduksiyon acisi ve goreli rotasyonlar, kullanicinin govdesine daha bagli
ve daha anlamli bicimde uretilir.

\subsection{Kalibrasyon Tabanli Goreli Rotasyonlar}

Sistem, onkol tracker'i, ust kol tracker'i ve el bilegi pozu arasindaki iliskiyi
mutlak dunya rotasyonlariyla degil, kalibrasyonla sifirlanmis goreli rotasyonlarla
hesaplar. Temel iliskiler sunlardir:

\begin{align}
    Q_{\text{wrist,rel}} &= Q_{\text{forearmTracker}}^{-1} Q_{\text{wrist}}
    \label{eq:wrist_rel} \\
    Q_{\text{elbow,rel}} &= Q_{\text{upperArmTracker}}^{-1} Q_{\text{forearmTracker}}
    \label{eq:elbow_rel} \\
    Q_{\text{shoulder,rel}} &= Q_{\text{ref}}^{-1} Q_{\text{upperArmTracker}}
    \label{eq:shoulder_rel}
\end{align}

Kalibrasyon aninda olculen notr rotasyonlar $Q_{\text{neutral}}$ olarak
kaydedilir ve calisma zamaninda:

\begin{equation}
    \Delta Q = Q_{\text{neutral}}^{-1} Q_{\text{rel}}
    \label{eq:delta_q}
\end{equation}

seklinde sifir referansli hareket elde edilir. Bu, kullanicinin tracker'i kola
tam ideal acida takmamis oldugu durumlarda bile sistemin ayni anatomik pozu
tekrar tekrar benzer sekilde yorumlamasini saglar.

\subsection{Eklem Pozlarinin Yeniden Kurulmasi}

El bilegi pozu OpenXR el izleme verisinden, dirsek ve omuz pozlari ise
anatomik uzunluklar ve tracker eksenleri kullanilarak turetilir. Forearm
yonu $\hat{f}$ ve upper-arm yonu $\hat{u}$ olmak uzere:

\begin{align}
    \mathbf{p}_{\text{elbow}} &= \mathbf{p}_{\text{wrist}} + \hat{f}\,L_f
    \label{eq:elbow_pos} \\
    \mathbf{p}_{\text{shoulder}} &= \mathbf{p}_{\text{elbow}} + \hat{u}\,L_u
    \label{eq:shoulder_pos}
\end{align}

Burada $L_f$ onkol, $L_u$ ise ust kol uzunlugudur. Bu geometri sayesinde
sistem yalnizca aci okumakla kalmaz; ayni zamanda omuz-dirsek-bilek zincirini
uzayda yeniden kurarak gorsellestirme ve avatar yansitma katmanlarini besler.

\subsection{Pose Follower ve Retargeter Katmani}

Kararlilik yalnizca veri uretiminde degil, verinin gorsel modele aktariminda da
gereklidir. Bu nedenle iki farkli sunum katmani eklenmistir:

\begin{table}[H]
\centering
\caption{Ust ekstremite gorsellestirme katmanlari}
\label{tab:ustextremite_katman}
\renewcommand{\arraystretch}{1.22}
\begin{tabularx}{\textwidth}{lX}
\toprule
\textbf{Bilesen} & \textbf{Islev} \\
\midrule
\texttt{UpperLimbAvatarPoseFollower} & Dunya uzayinda hesaplanan omuz-dirsek-bilek pozlarini dogrudan preview rig'e uygular \\
\texttt{UpperLimbAvatarRetargeter} & Goreli omuz/dirsek/bilek rotasyonlarini farkli bir avatar kol hiyerarsisine aktarir \\
\texttt{UpperLimbKinematicsSource} & Bu iki bilesen icin ortak veri sozlesmesini tanimlar \\
\bottomrule
\end{tabularx}
\end{table}

\texttt{UpperLimbAvatarPoseFollower}, takipten gelen eklem pozlarini
istenirse bir \texttt{previewRoot} altina yeniden esitler. Bu esitleme,
kaynak omuz referans pozu ile hedef preview koku arasinda bir donusum
matrisi tanimlayarak yapilir:

\begin{equation}
    M_{\text{map}} = T_{\text{preview}}\,T_{\text{sourceRef}}^{-1}
    \label{eq:preview_map}
\end{equation}

Bu sayede ayni takip verisi, kullanicidan farkli bir konumda duran ghost veya
uzak avatar uzerinde de bozulmadan gosterilebilir.\
\texttt{UpperLimbAvatarRetargeter} ise segment uzunluklarini degil, dogrudan
kalibrasyon-sonrasi goreli rotasyonlari uyguladigi icin farkli avatar oranlarinda
bile daha kararlı bir gorunum sunar.

\subsection{Donanimsiz Test ve Stabilize Up-Vector Mantigi}

Canli tracker verisi olmadan gelistirme yapabilmek icin
\texttt{UpperLimbKinematicsTestSource} yazilmistir. Bu bilesen,
sentetik omuz sallanimi ve dirsek fleksiyonu uretir; boylece pose follower,
retargeter ve gorsellestirme katmani editor icinde test edilebilir.

Bu test kaynaginda ileri yon vektoru ile uyumlu bir ``up'' vektoru elde etmek
icin projeksiyon ve onceki kare bilgisi birlikte kullanilir. $\mathbf{u}_{t-1}$
onceki kareden gelen onbellekli up vektoru olmak uzere:

\begin{equation}
    \mathbf{u}_t = \mathrm{normalize}\left(\mathrm{Slerp}(\mathbf{u}_{t-1},\;\Pi_{\perp \mathbf{f}}(\mathbf{u}_{\text{aday}}),\;0.5)\right)
    \label{eq:stable_up}
\end{equation}

Bu yaklasim, ani eksen ters donmelerini ve segment mesh'lerinde gorulebilen
titremeyi azaltir. Dolayisiyla hafta~4 calismasi sadece ``poz hesaplama''
degil, ayni zamanda \textbf{stabil sunum} katmanini da kurmustur.

% ================================================================
\section{Bitirme Tezi Taslaginin Hazirlanmasi}
\label{sec:thesis}

Hafta~4'un ucuncu buyuk ciktisi, projenin daginik teknik notlarini tek bir
akademik omurgada toplayan bitirme tezi taslagidir. Hazirlanan
\texttt{BitirmeTezi\_Taslak.tex} dosyasi, universitenin tez formatina uygun
kapak, onsoz, etik beyanname, ozet, simgeler dizini ve bolumlendirme duzenine
sahiptir.

\subsection{Taslagin Akademik Islevi}

Bu belge yalnizca ``ileride yazilacak tez icin bos bir kabuk'' degildir.
Mevcut durumda bile;

\begin{itemize}[nosep]
    \item proje amacini,
    \item teknik mimariyi,
    \item alt ekstremite ve full-body alt sistemlerini,
    \item oyunlastirilmis gorev motorunu,
    \item temel ciktilari, sinirliliklari ve gelecek calismalari
\end{itemize}

tek bir savunulabilir akista toplamaktadir.

\subsection{Taslakta Yer Alan Ana Bolumler}

\begin{table}[H]
\centering
\caption{Bitirme tezi taslaginin ana bolumleri}
\label{tab:tez_bolumleri}
\renewcommand{\arraystretch}{1.22}
\begin{tabularx}{\textwidth}{lX}
\toprule
\textbf{Bolum} & \textbf{Icerik} \\
\midrule
On kisim & Kapaklar, onsoz, etik beyanname, ozet, sekiller/tablolar dizini, kisaltmalar \\
\texttt{Genel Bilgiler} & Giris, problem tanimi, motivasyon, literatur ve teknik arka plan \\
\texttt{Yapilan Calismalar} & Sistem hedefleri, yazilim mimarisi, alt sistemler, gorev motoru ve gelisim asamalari \\
\texttt{Bulgular ve Tartisma} & Elde edilen teknik ciktilar, katkilar ve sinirliliklar \\
\texttt{Sonuclar} & Genel degerlendirme ve gelecek calismalar \\
Ekler & Gorev protokolu, acik birakilan alanlar, deneysel dogrulama plani \\
\bottomrule
\end{tabularx}
\end{table}

Bu yapinin onemi sudur: hafta hafta gelistirilen teknik bilesenler artik
ayri notlar halinde kalmamakta, tek bir akademik anlatinin parcasi haline
gelmektedir. Ozellikle \texttt{Yazilim mimarisi}, \texttt{Alt ekstremite ROM alt
sistemi}, \texttt{Full-body avatar ve IK alt sistemi} ve
\texttt{Oyunlastirilmis tarama ve gorev motoru} basliklari, proje kod tabaninin
tez duzeyinde savunulabilir bir sekilde aktarilmasini saglamaktadir.

\subsection{Hafta 4 Acisindan Anlami}

Hafta~4'te yazilan tez taslagi iki nedenle kritik bir esik olusturur:

\begin{enumerate}[nosep]
    \item Proje bundan sonra yalnizca teknik olarak calisan bir prototip degil,
          ayni zamanda \textbf{raporlanabilir bir bitirme calismasi} haline gelmistir.
    \item Sonraki haftalarda uretilecek deneysel sonuc, ekran goruntusu ve
          dogrulama verileri icin hazir bir belge iskeleti olusmustur.
\end{enumerate}

Kisacasi, bu hafta yalnizca yeni bir ozellik eklenmemis; projenin akademik
teslim yolu da somutlastirilmistir.

% ================================================================
\section{Hafta 4 ciktilari ve Degerlendirme}
\label{sec:sonuc}

Hafta~4 sonunda proje uc farkli eksende ilerlemistir:

\begin{table}[H]
\centering
\caption{Hafta 4 somut ciktilari}
\label{tab:week4_outputs}
\renewcommand{\arraystretch}{1.22}
\begin{tabularx}{\textwidth}{lX}
\toprule
\textbf{Cikti} & \textbf{Saglanan fayda} \\
\midrule
Ghost avatar demo sistemi & Kullaniciya gorevi metinle degil hareketle gosterebilme \\
Pose/sequence asset modeli & Demo hareketlerini kod degistirmeden guncelleyebilme \\
Editor otomasyonu & Sahne kurulumu ve gorev asset uretimini hizlandirma \\
Ust ekstremite referans cercevesi & Omuz/dirsek/bilek verisini daha anlamli ve daha stabil yorumlama \\
Pose follower + retargeter & Takip verisini farkli avatar ve preview rig'lerde guvenli gosterme \\
Tez taslagi & Projenin akademik teslimine hazir bir yapi kazandirma \\
\bottomrule
\end{tabularx}
\end{table}

Bu haftanin genel sonucu, sistemin yalnizca hesap yapan bir XR prototipi olmaktan
cikarak; kullaniciyi yonlendiren, cikan veriyi daha tutarli gosteren ve
akademik olarak savunulabilir bicimde kayda gecirilen bir calisma haline
gelmesidir.

Bir sonraki mantikli adimlar sunlardir:

\begin{enumerate}[nosep]
    \item Ghost avatar gosterimlerinin gercek kullanici oturumlarinda anlasilabilirlik acisindan gozlenmesi,
    \item Ust ekstremite hattinin canli tracker oturumlarinda tekrar edilebilirlik testleriyle dogrulanmasi,
    \item Tez taslagina nicel deney sonuclari, ekran goruntuleri ve kaynakca son halinin eklenmesi.
\end{enumerate}

\end{document}